\section{Introduction}
\label{sec:introduction}

Digital media cross local trust boundaries through email, web download, cloud
synchronization, direct import, and messaging. Upstream controls can establish
channel authenticity, access rights, or confidentiality, but they do not prove
that the received octet string is admissible to a local parser or decoder. It
may still preserve privacy-bearing metadata, amplify resource demand, exploit a
decoder, or admit conflicting parser interpretations. End-to-end encrypted
messaging is one instance because a blind relay cannot inspect plaintext
without changing the confidentiality model
\cite{unger2015,cohngordon2020,abelson2024}. The underlying file-admission
problem is not specific to that transport.

A suffix and caller-supplied media type are declarations, a magic-byte match
may cover only a prefix, and a signature scan cannot establish whole-object
grammar or removal of excluded structures. A successful transcoder exit proves
only that one tool emitted bytes. Each mechanism supplies evidence, but none
alone authorizes the emitted object for an endpoint filesystem or renderer.

Content disarm and reconstruction methods address part of this problem by
replacing input with a reduced representation. Prior work establishes the
importance of parser agreement, complete consumption, generated parsers, and
canonical languages \cite{samuel2012,jana2012,momot2016,nail2014,
everparse2019,ipg2023,darpaSafeDocs}. The unresolved systems question addressed
here is how route selection, reconstruction strength, complete-output
authorization, and filesystem publication compose at one local trust boundary.

This study develops a transport-agnostic method for policy-indexed canonical
media reconstruction and fail-closed delivery. \emph{Delivery} is any exposure
of authorized bytes through the component interface; \emph{publication} is the
narrower non-overwriting filesystem transition. The method converges portable
name, supplied type, and bounded byte evidence to one source profile, enforces
the policy-required reconstruction level, reparses the complete candidate on a
code path separate from construction, applies rebuilt-output scanning and
resource bounds, and permits publication only after a clean verdict. Every
other outcome exposes no deliverable bytes and leaves any pre-existing target
unchanged.

The analytical claim is conditional: under registered-parser, semantic-
construction, and filesystem assumptions, accepted output belongs to one
policy-indexed language and a non-deliverable result cannot publish. Empirical
conformance is falsified when a manifest record violates any prespecified
action, reason, state, reconstruction, closure, resource, or fault oracle.
Fidelity remains a separate usability condition and cannot authorize an
invalid structure.

The novelty is the authorization relation: successful transformation is not
authorization, reconstruction strength is policy checked, and publication is a
separate guarded transition. Four linked contributions follow. \emph{C1
(method)} composes unique routing, a registered source--level--output mapping,
reconstruction sufficiency, complete candidate validation, rebuilt-output
scanning, resource bounds, and clean-only delivery. \emph{C2 (analysis)} states
the assumption-to-claim boundary for conditional canonical delivery and fixed-
witness post-reconstruction monotonicity. \emph{C3 (system)} realizes the eight
auditable obligations F1--F8 with typed failures, validation separate from
construction, and non-overwriting publication. \emph{C4 (evaluation protocol)}
binds each obligation to record-level falsifiers, retains negative and
incomplete runs, and withholds controlled or release claims whenever a
scientific gate is not passed. C4 is a protocol contribution, not a completed
population-level empirical result.

Research question 1 (RQ1) asks when heterogeneous endpoint evidence selects
exactly one admissible route. RQ2 asks which conditions are sufficient to
restrict canonical delivery to clean verdicts at each reconstruction level.
RQ3 asks whether the implementation preserves those conditions on a frozen
corpus under malformed
and ambiguous carriers, decoder and publication faults, cancellation, and fixed
resource limits, and what fidelity and latency costs accompany the observed
paths. The first two questions receive analytical answers. RQ3 is only
partially exercised by the bounded finite-corpus experiment and remains
unanswered at its declared controlled and population scope.

The scope is measurable carrier structure. Malware family and operational role
require behavioral evidence unavailable from media bytes alone. The method does
not remove harmful decoded meaning, prove absence of steganography, compensate
for a compromised endpoint, or make an unsafe decoder trustworthy. For frozen
profiles and explicit assumptions, it restricts publication to candidates
validated separately from construction and produced at or above the required
reconstruction level.
